\subsubsection{有限 RH：s-变量 zeta 的临界带极点与 primitive 的平方根振荡}\label{app:real-input-40-finite-rh}
把 $z$-变量换为 Dirichlet 坐标
\[
z=\lambda^{-s},\qquad \widehat{\zeta}(s):=\zeta_M(\lambda^{-s})=\frac{1}{\det(I-\lambda^{-s}M)}.
\]
则 $\widehat{\zeta}(s)$ 的极点由 $M$ 的非零特征值 $\mu$ 决定：当且仅当
\[
\lambda^{-s}\mu=1
\]
时出现极点（对每个 $\mu$ 及其对数分支给出一列极点）。

\paragraph{谱—极点—误差三角形：有限维 \texorpdfstring{$\zeta$}{zeta} 的通用字典}
对任意有限维核（矩阵）$M$，Dirichlet 化后的 $\widehat{\zeta}(s)$ 把“谱模长”逐点翻译为“临界带极点的位置”，并进一步把“非平凡谱半径”翻译为 primitive 计数的误差尺度。
该字典是有限维模型中最接近 \(\zeta\)/\(L\) 语言的部分：它把“零点/极点—显式公式—误差项”的分析数论三角形压缩为一个纯线性代数三角形。

\begin{proposition}[谱 \texorpdfstring{$\leftrightarrow$}{<->} 极点：临界带坐标与 toy-RH 判别]\label{prop:finite-rh-triangle-dict}
设 $M$ 为有限维矩阵，$\lambda=\rho(M)>1$，并定义
\[
\widehat{\zeta}(s)=\zeta_M(\lambda^{-s})=\frac{1}{\det(I-\lambda^{-s}M)}.
\]
则对任意非零特征值 $\mu\neq 0$ 与任意 $k\in\ZZ$，$\widehat{\zeta}(s)$ 的极点由
\[
\lambda^{-s}\mu=1
\]
给出，等价于
\[
s=\frac{\log\mu+2\pi i k}{\log\lambda},
\]
其中 $\log\mu$ 取任意复对数分支。于是极点实部满足
\[
\Re(s)=\frac{\log|\mu|}{\log\lambda}.
\]
因此在开带 $0<\Re(s)<1$ 中出现非平凡极点当且仅当存在特征值满足 $1<|\mu|<\lambda$。

进一步，若把“有限维 toy-RH”定义为：开带 $0<\Re(s)<1$ 中的全部非平凡极点都落在 $\Re(s)=\tfrac12$，则该条件等价于
\[
(\forall\,\mu\ \text{满足}\ 1<|\mu|<\lambda)\qquad |\mu|=\sqrt{\lambda}.
\]
更弱地，若满足 Ramanujan 半径界
\[
\max_{\mu\neq\lambda}|\mu|\ \le\ \sqrt{\lambda},
\]
则右半开带 $\tfrac12<\Re(s)<1$ 内无极点；若存在 $|\mu|=\sqrt{\lambda}$，则对应极点列恰落在 $\Re(s)=\tfrac12$ 上。
\end{proposition}
\leanverified{paper\_finite\_rh\_triangle\_dict}

\begin{corollary}[kernel RH \texorpdfstring{$\Longleftrightarrow$}{<->} 半条带无极点]\label{cor:finite-rh-halfstrip-nopoles}
\leanverified{paper\_finite\_rh\_halfstrip\_nopoles}
沿用命题 \ref{prop:finite-rh-triangle-dict} 的记号。把 kernel RH（Ramanujan 半径界）表述为
\[
\max_{\mu\neq \lambda}|\mu|\le \sqrt{\lambda}.
\]
则该条件等价于：\(\widehat{\zeta}(s)\) 在开半条带
\[
\tfrac12<\Re(s)<1
\]
内无极点。并且该条件达到饱和（存在 \(|\mu|=\sqrt{\lambda}\)）当且仅当 \(\widehat{\zeta}(s)\) 在临界线 \(\Re(s)=\tfrac12\) 上出现非平凡极点列。
\leanverified{paper\_finite\_rh\_halfstrip\_nopoles}
\end{corollary}
\begin{proof}
由命题 \ref{prop:finite-rh-triangle-dict} 的极点实部恒等式 \(\Re(s)=\log|\mu|/\log\lambda\) 可知：存在某个 \(\mu\neq 0\) 使其极点落在 \(\tfrac12<\Re(s)<1\) 当且仅当 \(\sqrt{\lambda}<|\mu|<\lambda\)。
而 \(\lambda\) 为谱半径使得 \(|\mu|\le \lambda\) 对所有特征值成立，因此“半条带无极点”等价于对所有 \(\mu\neq\lambda\) 有 \(|\mu|\le \sqrt{\lambda}\)，即 kernel RH。
最后一段同理由 \(|\mu|=\sqrt{\lambda}\iff \Re(s)=\tfrac12\) 得到。证毕。
\end{proof}

\begin{remark}[本核的达界位置：唯一临界模长 \texorpdfstring{$|\mu|=\sqrt{\lambda}$}{|mu|=sqrt(lambda)}]\label{rem:finite-rh-40-ramanujan}
在真实输入 40 状态核上，除 Perron 根 $\lambda=\varphi^2$ 外，唯一满足 $|\mu|=\sqrt{\lambda}$ 的特征值为 $-\sqrt{\lambda}=-\varphi$（命题 \ref{prop:real-input-40-essential-20} 与此处命题 \ref{prop:finite-rh-40} 的谱分解一致）。
因此该核达到 toy-RH 的“单频临界振荡”情形：临界带中的全部非平凡极点由 $-\sqrt{\lambda}$ 单独产生。
若令 \(y:=\mu/\sqrt{\lambda}\) 并取 \(S=y+y^{-1}\)，则该达界点对应 \(y=-1\) 与 \(S=-2\)，与 \S\ref{subsec:xi-from-self-dual} 的完成化压缩坐标一致（见注 \ref{rem:xi-joukowski-toy-rh-alignment}）。
\end{remark}

\begin{remark}[临界线晶格的来源：一个隐藏的 \texorpdfstring{$(-A_0)$}{(-A0)} 因子]\label{rem:finite-rh-40-hidden-minusA0}
命题 \ref{prop:real-input-40-A0-tensor-hidden} 给出了更细的三因子分解
\[
\det(I-zM)
=(1-z)^2(1+z)\,
\det\!\bigl(I-z(A_0\otimes A_0)\bigr)\,
\det\!\bigl(I-z(-A_0)\bigr).
\]
其中 \(\det(I-z(-A_0))=\det(I+zA_0)=1+z-z^2\) 的根为
\(
z=\frac{1}{\varphi},\ -\varphi
\)，
从而在 \(M\) 的非零谱中强制出现特征值 \(-\varphi=-\sqrt{\lambda}\)。
因此 toy-RH 的“唯一达界”并非来自 \(40\) 状态复杂结构的偶然共振，而是由一个被嵌入进核的二维黄金均值因子 \((-A_0)\) 所必然产生。
进一步地，由块结构诱导的常纤维子空间 \(U\) 与垂直子空间 \(U^\perp\) 给出“骨架旧谱 / 纵向残差新谱”的两通道分解（见命题 \ref{prop:real-input-40-skeleton-factor} 与注 \ref{rem:real-input-40-deterministic-lift}）。因此 toy-(G)RH 的平方根界可被理解为：把全部非平凡谱约束到纵向残差通道中并施加相应的谱半径界。
\end{remark}

\begin{remark}[零点—振荡字典：临界圆谱 \texorpdfstring{$\Rightarrow$}{=>} 临界线梳齿与多频平方根项]\label{rem:finite-rh-zero-oscillation-dict}
命题 \ref{prop:finite-rh-triangle-dict} 给出了“模长 \(|\mu|\)”决定极点所在竖直直线的规则；若进一步把临界圆上的\textbf{相位}写出来，则它直接决定临界线上的“频率梳齿”，并在周期迹与 primitive 轨道计数中产生显式振荡项。
\par
设 \(\lambda=\rho(M)>1\)，并假设在临界圆 \(|\mu|=\sqrt{\lambda}\) 上不存在 Jordan 块（例如该模长部分半单）。记临界圆上的全部达界特征值为
\[
\mu_j=\sqrt{\lambda}\,e^{i\alpha_j}\qquad(j=1,\dots,J),
\]
按代数重数计；并设其余特征值的模长严格小于 \(\sqrt{\lambda}\)，令
\[
\rho_{\mathrm{sub}}:=\max\{|\mu|:\ \mu\in\mathrm{spec}(M),\ |\mu|<\sqrt{\lambda}\}<\sqrt{\lambda}.
\]
则：
\begin{enumerate}
  \item \textbf{（临界线极点梳齿）} 对每个 \(j\) 与每个 \(k\in\ZZ\)，由 \(\lambda^{-s}\mu_j=1\) 得到一条临界线极点列
  \[
  s=\frac12+\frac{i(\alpha_j+2\pi k)}{\log\lambda}.
  \]
  \item \textbf{（周期迹的多频平方根项）} 由谱分解（或对角化）有
  \[
  \Tr(M^n)=\lambda^n+\lambda^{n/2}\sum_{j=1}^{J}e^{in\alpha_j}+O\!\bigl(\rho_{\mathrm{sub}}^{\,n}\bigr).
  \]
  \item \textbf{（primitive 的振荡项）} 由 Witt/Möbius 反演
  \(p_n=\frac{1}{n}\sum_{d\mid n}\mu(d)\Tr(M^{n/d})\)
  得到
  \[
  p_n=\frac{\lambda^n}{n}+\frac{\lambda^{n/2}}{n}\sum_{j=1}^{J}e^{in\alpha_j}
  +O\!\left(\frac{\rho_{\mathrm{sub}}^{\,n}}{n}\right),
  \]
  其中对偶数 \(n\) 的细节（例如 \(\alpha_j=\pi\) 时的奇偶消项）可由命题 \ref{prop:finite-rh-parity-general} 进一步精化。
\end{enumerate}
在真实输入 40 状态核上，唯一达界相位为 \(\alpha=\pi\)（即 \(\mu=-\sqrt{\lambda}\)），因此临界线梳齿退化为单一等间距晶格，并且平方根项退化为 \(({-}1)^n\lambda^{n/2}\) 的单频振荡（推论 \ref{cor:finite-rh-explicit-formula}）。
\end{remark}

\paragraph{相位增殖：保持 RH-sharp 的循环提升}
\(\zeta(s)\) 的困难部分不仅是“实部落在 \(\tfrac12\)”这一条临界线，而是临界线上存在大量、且虚部结构复杂的零点族。
在有限维核模型中，可以在\textbf{保持 Ramanujan/RH-sharp 半径界}的同时，把“单一相位”的临界极点列增殖为任意细的相位梳齿；这为后续把 toy-RH 推向 toy-GRH（通过群扩张/Artin 分解）提供了一个最小可审计机制。

\begin{proposition}[循环相位提升：谱相位细分且保持 Ramanujan 半径界]\label{prop:finite-rh-phase-lift}
\leanverified{paper\_finite\_rh\_phase\_lift}
设 $M$ 为有限维矩阵，$\lambda=\rho(M)>1$。取整数 $q\ge 2$，令 $P_q$ 为 $q$-循环置换矩阵（$P_q e_j=e_{j+1}$，指标按 $q$ 取模），并定义循环提升矩阵
\[
M^{\langle q\rangle}:=M\otimes P_q.
\]
则：
\begin{enumerate}
  \item \textbf{（谱相位细分）} 设 $\omega_q=e^{2\pi i/q}$。有谱恒等式
  \[
  \mathrm{spec}(M^{\langle q\rangle})
  =\{\mu\,\omega_q^k:\ \mu\in \mathrm{spec}(M),\ k=0,1,\dots,q-1\},
  \]
  按代数重数计。
  \item \textbf{（谱半径保持）} $\rho(M^{\langle q\rangle})=\rho(M)=\lambda$。
  \item \textbf{（Ramanujan 界保持）} 若 $M$ 满足 $\max_{\mu\neq \lambda}|\mu|\le \sqrt{\lambda}$，则 $M^{\langle q\rangle}$ 亦满足同一半径界。
  \item \textbf{（临界线极点梳齿）} 令
  \(
  \widehat{\zeta}^{\langle q\rangle}(s):=\det(I-\lambda^{-s}M^{\langle q\rangle})^{-1}
  \)。
  若 $\mu=\sqrt{\lambda}e^{i\theta}$ 是 $M$ 的某个达界特征值，则它在 $M^{\langle q\rangle}$ 上生成 $q$ 条互相平移的临界线极点列：
  \[
  s=\frac12+\frac{i}{\log\lambda}\Bigl(\theta+2\pi\Bigl(\frac{k}{q}+m\Bigr)\Bigr),
  \qquad k\in\{0,\dots,q-1\},\ m\in\ZZ.
  \]
\end{enumerate}
当 $q\to\infty$ 时，上述极点在每个基本周期窗 $[0,2\pi/\log\lambda)$ 内的虚部位置变得任意稠密。
\end{proposition}

\begin{proof}
置换矩阵 $P_q$ 可被离散傅里叶变换对角化，其特征值为 $\omega_q^k$（$0\le k\le q-1$）。
对任意 $M$ 的特征对 $(\mu,v)$ 与任意 $P_q$ 的特征对 $(\omega_q^k,w_k)$，张量向量 $v\otimes w_k$ 满足
\(
(M\otimes P_q)(v\otimes w_k)=(Mv)\otimes(P_qw_k)=\mu\,\omega_q^k (v\otimes w_k)
\)，从而得到谱相位细分公式；反向包含由维数计数或 Kronecker 谱性质得到。
谱半径与 Ramanujan 界的保持性来自 $|\omega_q^k|=1$。
最后极点公式由命题 \ref{prop:finite-rh-triangle-dict} 对 $M^{\langle q\rangle}$ 逐个特征值应用并取达界模长 $|\mu|=\sqrt{\lambda}$ 得到。证毕。
\end{proof}

\begin{remark}[应用到真实输入 $40$ 状态核：单频 \texorpdfstring{$\pi$}{pi} 振荡的细分]\label{rem:finite-rh-40-phase-lift}
对真实输入 $40$ 状态核，唯一达界特征值为 $-\sqrt{\lambda}=-\varphi=\sqrt{\lambda}e^{i\pi}$（注 \ref{rem:finite-rh-40-ramanujan}）。
因此循环相位提升把原来的单一临界列
\(
s=\tfrac12+\frac{(2m+1)\pi i}{\log\lambda}
\)
细分为
\[
s=\frac12+\frac{i}{\log\lambda}\Bigl(\pi+2\pi\Bigl(\frac{k}{q}+m\Bigr)\Bigr),
\qquad k=0,\dots,q-1,\ m\in\ZZ,
\]
在保持 RH-sharp 的同时把临界线上的虚部结构从“单栅栏”提升为“$q$-梳齿”。
\end{remark}

\begin{proposition}[循环提升的 Artin 因子化（阿贝尔角色通道的显式分解）]\label{prop:finite-rh-phase-lift-artin}
\leanverified{paper\_finite\_rh\_phase\_lift\_artin}
沿用命题 \ref{prop:finite-rh-phase-lift} 的记号。令
\(
\omega_q=e^{2\pi i/q}
\)，并把循环群 \(G=\ZZ/q\ZZ\) 的角色记为 \(\chi_k(j)=\omega_q^{kj}\)（\(k\in\{0,\dots,q-1\}\)）。
则对任意复参数 \(z\) 有行列式分解
\[
\det\!\bigl(I-zM^{\langle q\rangle}\bigr)
=\prod_{k=0}^{q-1}\det\!\bigl(I-z\,\omega_q^k M\bigr).
\]
等价地，循环提升的 \(\zeta\)-函数满足
\[
\zeta_{M^{\langle q\rangle}}(z)
=\prod_{k=0}^{q-1}\zeta_{M}(z\omega_q^k),
\qquad
\zeta_M(z):=\det(I-zM)^{-1}.
\]
因此“相位增殖”不仅是谱点集合的旋转细分，更是一个可审计的\textbf{阿贝尔 Artin 分解}：每个角色通道 \(\chi_k\) 对应一个“扭曲因子” \(\det(I-z\,\omega_q^k M)\)。
\end{proposition}

\begin{proof}
取离散傅里叶矩阵 \(F\) 使得 \(F^{-1}P_qF=\mathrm{diag}(\omega_q^0,\dots,\omega_q^{q-1})\)。
则
\[
I-z(M\otimes P_q)
\sim
I-z\bigl(M\otimes \mathrm{diag}(\omega_q^0,\dots,\omega_q^{q-1})\bigr)
\]
在相似意义下等价于一个按 \(k\) 分块的块对角矩阵，其第 \(k\) 个块为 \(I-z\,\omega_q^k M\)。对块对角矩阵取行列式并相乘即得结论。证毕。
\end{proof}

\begin{corollary}[\texorpdfstring{$q$}{q} 步时间重标度恒等式：循环提升等价于 \texorpdfstring{$z\mapsto z^q$}{z->z^q} 的拉回]\label{cor:finite-rh-phase-lift-time-rescale}
\leanverified{paper\_finite\_rh\_phase\_lift\_time\_rescale}
沿用命题 \ref{prop:finite-rh-phase-lift} 的记号，令 \(M^{\langle q\rangle}=M\otimes P_q\)。
则对任意复参数 \(z\) 有恒等式
\[
\det\!\bigl(I-zM^{\langle q\rangle}\bigr)=\det\!\bigl(I-z^q M^q\bigr).
\]
等价地，
\[
\zeta_{M^{\langle q\rangle}}(z)=\zeta_{M^q}(z^q).
\]
在 Dirichlet 坐标 \(z=\lambda^{-s}\) 下，上式可改写为
\[
\widehat{\zeta}^{\langle q\rangle}(s)
=\zeta_{M^{\langle q\rangle}}(\lambda^{-s})
=\zeta_{M^q}(\lambda^{-qs}).
\]
\end{corollary}

\begin{proof}
由命题 \ref{prop:finite-rh-phase-lift-artin} 的行列式分解与恒等式
\(
\prod_{k=0}^{q-1}(1-x\omega_q^k)=1-x^q
\)，对任意 \(M\) 的特征值 \(\mu\) 有
\[
\prod_{k=0}^{q-1}(1-z\mu\omega_q^k)=1-(z\mu)^q=1-z^q\mu^q.
\]
对全部特征值按代数重数相乘得到
\(
\det(I-zM^{\langle q\rangle})=\prod_{\mu}(1-z^q\mu^q)=\det(I-z^qM^q)
\)。
取倒数即得 \(\zeta\)-恒等式；最后一式把 \(z=\lambda^{-s}\) 代入即可。证毕。
\end{proof}

\begin{corollary}[循环 lift 的迹序列整除筛选]\label{cor:finite-rh-phase-lift-trace-filter}
\leanverified{paper\_finite\_rh\_phase\_lift\_trace\_filter}
沿用命题 \ref{prop:finite-rh-phase-lift} 的记号，令 \(M^{\langle q\rangle}:=M\otimes P_q\)。
则对任意整数 \(n\ge 1\) 有
\[
\Tr\!\bigl((M^{\langle q\rangle})^{n}\bigr)
=\Tr(M^n)\,\Tr(P_q^n)
=
\begin{cases}
q\,\Tr(M^n), & q\mid n,\\[4pt]
0, & q\nmid n.
\end{cases}
\]
\end{corollary}
\begin{proof}
由 Kronecker 乘积性质 \((A\otimes B)^n=A^n\otimes B^n\) 与 \(\Tr(A\otimes B)=\Tr(A)\Tr(B)\) 得
\(\Tr((M^{\langle q\rangle})^n)=\Tr(M^n)\Tr(P_q^n)\)。
而 \(P_q^n\) 是 \(q\)-循环置换的 \(n\) 次幂，其迹等于不动点数：当且仅当 \(q\mid n\) 时为 \(q\)，否则为 \(0\)。
\end{proof}

\begin{remark}[机制含义：相位梳齿 \texorpdfstring{$\neq$}{!=} 相位混合]\label{rem:finite-rh-phase-lift-divisibility-filter}
推论 \ref{cor:finite-rh-phase-lift-trace-filter} 表明：循环 lift 在临界线极点侧实现“相位细分”，但在长度/迹侧同时引入了严格的整除门控——只有 \(n\equiv 0\pmod q\) 的长度层才承载 lift 引入的额外信息。
因此它更像产生“稀疏但幅度大的共振尖峰”，而不是在长度轴上处处出现的连续相位干涉；若要模拟显式公式中更细的振荡形态，需叠加更强的混合机制（例如多尺度并联或连续尺度混合）。
\end{remark}

\paragraph{连续相位极限：Jensen/Mahler 自由能与临界折点}\label{par:finite-rh-jensen-free-energy}
循环 lift 的 Artin 分解在“相位侧”提供了一个自然的极限机制：当相位梳齿足够细时，与其追踪单个极点的位置，不如改研究一个只依赖谱模长的宏观量（相位平均的 \(\log|\det|\)），它会在阈值
\(\sigma=\log|\mu|/\log\lambda\)
处产生不可导折点，特别地当存在达界模态 \(|\mu|=\sqrt{\lambda}\) 时，\(\sigma=\tfrac12\) 变成一个“相变点”。

\begin{proposition}[Jensen 平均恒等式：相位平均的 \texorpdfstring{$\log|\det|$}{log|det|} 只依赖谱模长]\label{prop:finite-rh-jensen-free-energy}
设 \(M\) 为有限维矩阵，\(\lambda=\rho(M)>1\)，其非零特征值（含重数）为 \(\{\mu_j\}_{j=1}^{J}\)。
对任意 \(\sigma\in\RR\) 定义
\[
G_M(\sigma):=\int_0^{2\pi}\log\Big|\det\bigl(I-\lambda^{-\sigma}e^{i\theta}M\bigr)\Big|\frac{d\theta}{2\pi}.
\]
则有闭式
\[
\boxed{\;
G_M(\sigma)=\sum_{j=1}^{J}\max\bigl(0,\ \log|\mu_j|-\sigma\log\lambda\bigr)\; }.
\]
并且对任意 \(q\ge 2\) 的循环提升 \(M^{\langle q\rangle}=M\otimes P_q\)，若定义（按维度归一化）
\[
G_{M^{\langle q\rangle}}(\sigma):=\frac{1}{q}\int_0^{2\pi}\log\Big|\det\bigl(I-\lambda^{-\sigma}e^{i\theta}M^{\langle q\rangle}\bigr)\Big|\frac{d\theta}{2\pi},
\]
则恒有 \(G_{M^{\langle q\rangle}}(\sigma)=G_M(\sigma)\)；因此 \(q\to\infty\) 的“连续相位极限”在该宏观量意义下是稳定的。
\end{proposition}
\begin{proof}
由特征值分解（只需按代数重数计）有
\[
\det\bigl(I-\lambda^{-\sigma}e^{i\theta}M\bigr)=\prod_{j=1}^{J}\bigl(1-\lambda^{-\sigma}e^{i\theta}\mu_j\bigr).
\]
对每一因子应用 Jensen 平均恒等式（对任意 \(a\in\CC\)）：
\[
\int_0^{2\pi}\log|1-ae^{i\theta}|\frac{d\theta}{2\pi}
=
\begin{cases}
0,&|a|\le 1,\\
\log|a|,&|a|>1,
\end{cases}
\]
即
\(
\int \log|1-ae^{i\theta}|= \max(0,\log|a|)
\)。
取 \(a=\lambda^{-\sigma}\mu_j\) 并求和，得到
\[
G_M(\sigma)=\sum_{j=1}^{J}\max\bigl(0,\ \log|\lambda^{-\sigma}\mu_j|\bigr)
=\sum_{j=1}^{J}\max\bigl(0,\ \log|\mu_j|-\sigma\log\lambda\bigr).
\]
对 lift 的结论由命题 \ref{prop:finite-rh-phase-lift} 的谱相位细分
\(\mathrm{spec}(M^{\langle q\rangle})=\{\mu_j\omega_q^k\}\)
立得：每个 \(\mu_j\) 在 lift 中出现 \(q\) 次且模长不变，于是
\[
G_{M^{\langle q\rangle}}(\sigma)=\frac{1}{q}\sum_{j=1}^J\sum_{k=0}^{q-1}\max\bigl(0,\log|\mu_j\omega_q^k|-\sigma\log\lambda\bigr)=\sum_{j=1}^J\max\bigl(0,\log|\mu_j|-\sigma\log\lambda\bigr).
\]
证毕。
\leanverified{paper\_finite\_rh\_jensen\_free\_energy}
\end{proof}

\begin{corollary}[真实输入 40 状态核的临界折点：\texorpdfstring{$\sigma=\tfrac12$}{sigma=1/2}]\label{cor:real-input-40-jensen-kink}
\leanverified{paper\_real\_input\_40\_jensen\_kink}
对真实输入 40 状态核，\(\lambda=\varphi^2\) 且其非零谱为
\(\{\lambda,\ -\sqrt{\lambda},\ 1^{(2)},\ (-1)^{(3)},\ \lambda^{-1},\ \lambda^{-1/2}\}\)。
因此在带 \(0\le \sigma\le 1\) 上，只有模长 \(>\!1\) 的两点 \(|\mu|=\lambda\) 与 \(|\mu|=\sqrt{\lambda}\) 贡献 \(G_M(\sigma)\)，从而
\[
G_M(\sigma)=
\begin{cases}
(1-\sigma)\log\lambda,&\tfrac12\le \sigma\le 1,\\[4pt]
\bigl(\tfrac32-2\sigma\bigr)\log\lambda,&0\le \sigma\le \tfrac12.
\end{cases}
\]
特别地 \(G_M\) 在 \(\sigma=\tfrac12\) 处出现不可导折点（斜率从 \(-\log\lambda\) 跳到 \(-2\log\lambda\)）。
\end{corollary}

\begin{remark}[高度--尺度校准：用 \texorpdfstring{$T=(2\pi)\lambda^q$}{T=(2pi)lambda^q} 对齐临界线梳齿密度]\label{rem:finite-rh-height-scale-calibration}
设 \(M\) 满足 RH-sharp，并令某个达界特征值写为 \(\mu=\sqrt{\lambda}\,e^{i\theta}\)。
由命题 \ref{prop:finite-rh-phase-lift}，循环提升 \(M^{\langle q\rangle}\) 在临界线 \(\Re(s)=\tfrac12\) 上生成的极点（toy-zero）虚部为
\[
\gamma_{k,m}=\frac{\theta+2\pi(\frac{k}{q}+m)}{\log\lambda},
\qquad k=0,1,\dots,q-1,\ m\in\ZZ,
\]
因此基本间距为
\[
\Delta\gamma_{\mathrm{toy}}=\frac{2\pi}{q\log\lambda}.
\]
\paragraph{标度校准}
若把“有效高度”定义为
\[
T:=(2\pi)\lambda^q
\quad\Longleftrightarrow\quad
\log\!\left(\frac{T}{2\pi}\right)=q\log\lambda,
\]
则上述间距可改写为
\[
\boxed{\ \Delta\gamma_{\mathrm{toy}}=\frac{2\pi}{\log(T/(2\pi))}\ }
\]
的对数尺度形式。
该结构与黎曼零点计数主项 \(N(T)\sim \frac{T}{2\pi}\log\frac{T}{2\pi}\) 导出的平均间距 \(2\pi/\log(T/(2\pi))\) 在形式上对齐；因此 lift 阶 \(q\) 可被解释为高度窗的对数尺度参数（\(q\) 增大对应更高高度窗的零点密度）。
\end{remark}

\paragraph{toy-explicit-formula 的反问题：零点工程（可计算的谱设计任务）}
给定目标虚部集合 \(\Gamma:=\{\gamma_1,\dots,\gamma_N\}\subset\RR\)（例如来自某个高度窗），固定一个满足 RH-sharp 的基核 \(M\) 及其 \(\lambda=\rho(M)\)。
利用注 \ref{rem:finite-rh-height-scale-calibration} 选取 lift 阶 \(q\) 使得 \(q\log\lambda\approx \log(T/(2\pi))\)（目标高度 \(T\) 的对数尺度）；随后用循环提升 \(M^{\langle q\rangle}\) 的临界线梳齿
\[
\gamma_{k,m}=\frac{\theta+2\pi(\frac{k}{q}+m)}{\log\lambda}
\]
对 \(\Gamma\) 进行拟合（选择整数对 \((k,m)\) 以最小化误差）。
\par
若单一 \((\lambda,q)\) 的严格晶格结构不足以复现目标的“非格点”局部统计，则可采用多尺度并联的最小机制：取若干满足 RH-sharp 的核 \(M^{(a)}\)（允许 \(\lambda_a\) 不同），并做块对角并合
\(
\bigoplus_a M^{(a)\langle q_a\rangle}
\)。
由于 \(\zeta\) 在块对角下乘性分解，其临界线极点集合为各通道梳齿的并；
当 \(\log\lambda_a\) 之间缺乏有理相关时，不同通道的周期 \(2\pi/\log\lambda_a\) 产生拍频，从而打破单一晶格的重复性。

\begin{proposition}[有限窗零点集的 \texorpdfstring{$\varepsilon$}{eps}-逼近（单尺度梳齿）]\label{prop:finite-rh-zero-engineering-eps}
\leanverified{paper\_finite\_rh\_zero\_engineering\_eps}
固定一个满足 RH-sharp 的基核 \(M\)，记 \(\lambda=\rho(M)>1\)，并取其某个达界特征值 \(\mu=\sqrt{\lambda}\,e^{i\theta}\)（\(\theta\in\RR\)）。
令 \(M^{\langle q\rangle}=M\otimes P_q\) 为循环提升（命题 \ref{prop:finite-rh-phase-lift}），并记由该达界模态在临界线 \(\Re(s)=\tfrac12\) 上生成的虚部梳齿为
\[
\gamma_{k,m}^{(q)}:=\frac{\theta+2\pi(\frac{k}{q}+m)}{\log\lambda},
\qquad k\in\{0,\dots,q-1\},\ m\in\ZZ.
\]
则对任意 \(\varepsilon>0\) 与任意有限集合 \(\Gamma=\{\gamma_1,\dots,\gamma_N\}\subset\RR\)，只要选取整数
\[
q\ \ge\ \frac{\pi}{\varepsilon\log\lambda},
\]
就存在整数对 \(\{(k_j,m_j)\}_{j=1}^N\) 使得对每个 \(j\) 都有
\[
\left|\gamma_j-\gamma_{k_j,m_j}^{(q)}\right|\le \varepsilon.
\]
\end{proposition}
\begin{proof}
对固定 \(q\)，上述集合满足
\[
\left\{\gamma_{k,m}^{(q)}:\ k=0,\dots,q-1,\ m\in\ZZ\right\}
=\left\{\frac{\theta+2\pi n/q}{\log\lambda}:\ n\in\ZZ\right\},
\]
其在实轴上构成等差晶格，步长为 \(\Delta\gamma_{\mathrm{toy}}=2\pi/(q\log\lambda)\)（注 \ref{rem:finite-rh-height-scale-calibration}）。
因此对任意 \(\gamma\in\RR\)，取
\[
\text{某个 }n\in\ZZ\text{ 使得 }\left|\gamma-\frac{\theta+2\pi n/q}{\log\lambda}\right|
\text{ 达到 }\min_{n'\in\ZZ}\left|\gamma-\frac{\theta+2\pi n'/q}{\log\lambda}\right|
\]
即可保证
\[
\min_{n'\in\ZZ}\left|\gamma-\frac{\theta+2\pi n'/q}{\log\lambda}\right|
\le \frac{1}{2}\Delta\gamma_{\mathrm{toy}}
=\frac{\pi}{q\log\lambda}.
\]
把所选 \(n\) 写成 \(n=k+qm\)（其中 \(k\in\{0,\dots,q-1\},\,m\in\ZZ\)），则得到对应的 \((k,m)\)。
对每个 \(\gamma_j\) 重复该选择，并令 \(q\ge \pi/(\varepsilon\log\lambda)\)，即得所示 \(\varepsilon\)-逼近。
\end{proof}

\begin{proposition}[拍频去晶格的必要条件：\texorpdfstring{$\log\lambda$}{loglambda} 的无理相关]\label{prop:finite-rh-beat-irrational-log}
设有两个 RH-sharp 通道 \((M^{(a)},\lambda_a)\)、\((M^{(b)},\lambda_b)\)，其中 \(\lambda_a,\lambda_b>1\)，并分别取循环提升阶 \(q_a,q_b\ge 1\)（允许 \(q_a=q_b=1\)）。
假设这两个通道在临界线 \(\Re(s)=\tfrac12\) 上各自产生的虚部点集包含至少一条完整等差梳齿列，其基本步长分别为
\[
\Delta_a=\frac{2\pi}{q_a\log\lambda_a},\qquad
\Delta_b=\frac{2\pi}{q_b\log\lambda_b}.
\]
若 \(\log\lambda_a/\log\lambda_b\in\QQ\)（等价地 \(\Delta_a/\Delta_b\in\QQ\)），则存在 \(\Delta>0\) 使得两通道的临界线虚部点集在 \(\Delta\) 下同时周期化；从而它们的并集仍是有限条算术晶格（等差竖线梳齿）的有限并，无法在任何有限尺度上摆脱周期锁相。
因此，若希望通过多通道并联获得\emph{无共同虚周期}的临界线相位结构（即“拍频去晶格”），必要地必须有
\[
\boxed{\ \frac{\log\lambda_a}{\log\lambda_b}\notin\QQ\ }.
\]
\end{proposition}
\begin{proof}
若 \(\log\lambda_a/\log\lambda_b=p/q\in\QQ\)（\(p,q\in\NN\)），则
\[
\frac{\Delta_a}{\Delta_b}
=\frac{q_b\log\lambda_b}{q_a\log\lambda_a}
=\frac{q_b q}{q_a p}\in\QQ.
\]
故存在整数 \(u,v\ge 1\) 使 \(u\Delta_a=v\Delta_b\)。取 \(\Delta:=u\Delta_a=v\Delta_b\)，则对任意形如 \(t_0+\Delta_a\ZZ\) 的梳齿列有
\[
(t_0+\Delta_a\ZZ)+\Delta=t_0+\Delta_a(\ZZ+u)=t_0+\Delta_a\ZZ,
\]
即其在平移 \(\Delta\) 下保持不变；同理对步长 \(\Delta_b\) 的梳齿列亦然。
由于每个通道的临界线点集是有限条梳齿列的并，故二者的并集在 \(\Delta\) 下周期化，从而仍属于有限条算术晶格的有限并。
最后一段为其逆否命题表述：若要排除任何非零共同虚周期，则必须有 \(\Delta_a/\Delta_b\notin\QQ\)，等价于 \(\log\lambda_a/\log\lambda_b\notin\QQ\)。
\end{proof}
\leanverified{paper\_finite\_rh\_beat\_irrational\_log}

\begin{proposition}[临界带中的非平凡极点都在 \texorpdfstring{$\Re(s)=\tfrac12$}{Re(s)=1/2}]\label{prop:finite-rh-40}
\leanverified{paper\_finite\_rh\_40}
设 $\lambda=\rho(M)=\varphi^2$。则在开带 $0<\Re(s)<1$ 中，$\widehat{\zeta}(s)$ 的全部极点均严格落在直线
\[
\Re(s)=\frac12.
\]
更具体地，唯一落在 $0<\Re(s)<1$ 的极点列来自特征值 $-\varphi=-\sqrt{\lambda}$，其位置为
\[
s=\frac12+\frac{(2k+1)\pi i}{\log\lambda},\qquad k\in\ZZ.
\]
\end{proposition}

\begin{proof}
由上文的谱分解，$M$ 的非零特征值集合为
\[
\{\lambda,\ -\sqrt{\lambda},\ 1^{(2)},\ (-1)^{(3)},\ \lambda^{-1},\ \lambda^{-1/2}\},
\]
其余 $31$ 个特征值为 $0$。对每个非零特征值 $\mu$，$\det(I-\lambda^{-s}M)=0$ 等价于 $\lambda^{-s}\mu=1$。

当 $\mu=\lambda$ 时得到主极点列 $s=1+2\pi i k/\log\lambda$，不落在开带 $0<\Re(s)<1$。
当 $\mu=-\sqrt{\lambda}$ 时，取复对数分支 $\log(-\sqrt{\lambda})=\tfrac12\log\lambda+i(2k+1)\pi$，
于是得到
\(
s=\tfrac12+\frac{(2k+1)\pi i}{\log\lambda}
\)，其满足 $0<\Re(s)<1$。
其余特征值满足 $|\mu|\le 1$，对应的极点满足 $\Re(s)\le 0$（落在临界带左侧边界或更左），因此在开带 $0<\Re(s)<1$ 内不产生极点。
\end{proof}

\begin{remark}[Dirichlet 变量与“临界线”表述]
令 $r:=\lambda^{1-s}$，则因子 $r+\varphi$ 对应的奇点满足
\[
r=-\varphi=\varphi e^{i(\pi+2\pi k)}.
\]
由 $\lambda=\varphi^2$ 得
\[
s=\frac12-\frac{(2k+1)\pi i}{\log\lambda},
\]
即在 $\Re(s)>0$ 内的非实奇点全部落在 $\Re(s)=1/2$ 上。
\end{remark}

\begin{proposition}[RH-sharp 的奇偶消项：Möbius 反演的一般机制]\label{prop:finite-rh-parity-general}
\leanverified{paper\_finite\_rh\_parity\_general}
设 \(\lambda>1\)，并给定一列实数 \(a_n\) 满足
\[
a_n=\lambda^n+\epsilon^n\lambda^{n/2}+O(\lambda^{n/3}),
\qquad \epsilon\in\{\pm 1\}.
\]
定义 primitive 轨道数（Witt/Möbius 反演）
\[
p_n:=\frac{1}{n}\sum_{d\mid n}\mu(d)\,a_{n/d}.
\]
则当 \(n\to\infty\) 时有：
\begin{enumerate}
  \item 若 \(n\) 为奇数，则
  \[
  p_n=\frac{\lambda^n}{n}+\frac{\epsilon^n\lambda^{n/2}}{n}+O\!\left(\frac{\lambda^{n/3}}{n}\right).
  \]
  \item 若 \(n\) 为偶数，则 \(d=1,2\) 两项在 \(\lambda^{n/2}\) 级别发生精确抵消，从而
  \[
  p_n=\frac{\lambda^n}{n}+O\!\left(\frac{\lambda^{n/3}}{n}\right).
  \]
\end{enumerate}
\end{proposition}
\begin{proof}
由定义
\[
np_n=\sum_{d\mid n}\mu(d)\,a_{n/d}.
\]
若 \(n\) 为奇数，则对所有 \(d\mid n\) 且 \(d\ge 2\) 有 \(d\ge 3\)，因此
\(
a_{n/d}=O(\lambda^{n/3})
\)。
于是
\[
np_n=a_n+\sum_{\substack{d\mid n\\ d\ge 3}}\mu(d)\,a_{n/d}
=\lambda^n+\epsilon^n\lambda^{n/2}+O(\lambda^{n/3}),
\]
得第一式。
\par
若 \(n\) 为偶数，则把 \(d=1,2\) 两项分离：
\[
np_n=a_n-a_{n/2}+\sum_{\substack{d\mid n\\ d\ge 3}}\mu(d)\,a_{n/d}.
\]
由假设，因 \(n\) 为偶数而 \(\epsilon^n=1\)，故
\[
a_n-a_{n/2}
\;=\;
(\lambda^n+\lambda^{n/2}+O(\lambda^{n/3}))-(\lambda^{n/2}+O(\lambda^{n/4}))
\;=\;
\lambda^n+O(\lambda^{n/3}).
\]
并且余项仍为 \(O(\lambda^{n/3})\)，合并得第二式。证毕。
\end{proof}

\begin{corollary}[显式公式形状：\texorpdfstring{$a_n$}{a\_n} 的平方根振荡与 primitive 的奇偶消项]\label{cor:finite-rh-explicit-formula}
\leanverified{paper\_finite\_rh\_explicit\_formula}
记 $a_n=\Tr(M^n)$。则
\[
a_n=\lambda^n+(-1)^n\lambda^{n/2}+O(1),
\]
并且 primitive 轨道数满足更精细的奇偶结构：
\[
p_n=\frac{\lambda^n}{n}-\frac{\lambda^{n/2}}{n}+O\!\left(\frac{\lambda^{n/3}}{n}\right)\qquad(n\ \text{为奇数}),
\]
\[
p_n=\frac{\lambda^n}{n}+O\!\left(\frac{\lambda^{n/3}}{n}\right)\qquad(n\ \text{为偶数}),
\]
其中偶数情形的 $\lambda^{n/2}$ 级项被 Möbius 反演中的 $d=1$ 与 $d=2$ 两项\textbf{精确抵消}。
\[
\boxed{
\;p_n=\frac{\lambda^n}{n}+\frac{\bigl((-1)^n-\mathbf{1}_{2\mid n}\bigr)\lambda^{n/2}}{n}
+O\!\left(\frac{\lambda^{n/3}}{n}\right)\; }
\]
并且具有可核验的“奇偶共振 signature”：
\[
\lim_{\substack{n\to\infty\\ n\ \mathrm{odd}}}\ \frac{n}{\lambda^{n/2}}\left(p_n-\frac{\lambda^n}{n}\right)=-1,
\qquad
\lim_{\substack{n\to\infty\\ n\ \mathrm{even}}}\ \frac{n}{\lambda^{n/2}}\left(p_n-\frac{\lambda^n}{n}\right)=0.
\]
\end{corollary}

\begin{remark}[“系数为 \(1\)”的平方根振荡：应在奇长层或去偶长干涉量上观测]\label{rem:finite-rh-unit-coefficient-observable}
推论 \ref{cor:finite-rh-explicit-formula} 中，周期迹层的平方根振荡项 \( (-1)^n\lambda^{n/2}\) 的\textbf{系数严格为 \(1\)}（其来源是唯一的达界特征值 \(-\sqrt{\lambda}\)）。
但在 primitive 层，Möbius 反演的 \(d=2\) 项会在偶长层对 \(\lambda^{n/2}\) 级主振荡产生\textbf{精确干涉抵消}，因此“系数为 \(1\)”应在如下两种等价的可检验观测量上读取：
\begin{itemize}
  \item \textbf{（奇长 signature）} 对奇数 \(n\)，有更强的归一化极限
  \[
  \boxed{\ \frac{n}{\lambda^{n/2}}\left(p_n-\frac{\lambda^n}{n}\right)=-1+O(\lambda^{-n/6})\ }.
  \]
  \item \textbf{（统一定义的共振指标）} 定义
  \[
  \mathsf{E}_n:=\frac{n}{\lambda^{n/2}}\left(p_n-\frac{\lambda^n}{n}\right),
  \]
  则 \(\mathsf{E}_n\to -1\)（\(n\) 奇）而 \(\mathsf{E}_n\to 0\)（\(n\) 偶），从而“单频、系数 \(1\)”的平方根振荡在 primitive 层以\textbf{奇偶门控}的方式出现。
\end{itemize}
这使 RH-sharp 的“平方根级误差达到”在数据侧具有一个极简审计接口：仅用 \(p_n\) 的奇长层（或等价的 \(\mathsf{E}_n\) 奇偶极限）即可直接观测到达界振荡；无需额外画极点列或引入外部拟合。
\end{remark}

\begin{proposition}[谱消去差分算子：有限谱的有限阶精确滤波]\label{prop:finite-rh-spectral-annihilation-filter}
\leanverified{paper\_finite\_rh\_spectral\_annihilation\_filter}
设 \(\{a_n\}_{n\ge 1}\) 是一个有限谱指数和：
\[
a_n=\sum_{j=1}^{J} m_j\,\mu_j^{\,n},
\]
其中 \(\mu_j\in\CC\)（按重数计）且 \(m_j\in\CC\)。
令移位算子 \(L\) 作用为 \((La)_n:=a_{n+1}\)。
对任意有限子集 \(S\subset\{\mu_1,\dots,\mu_J\}\) 定义多项式
\[
Q_S(x):=\prod_{\mu\in S}(x-\mu).
\]
则对所有 \(n\ge 1\) 有
\[
\boxed{\;
(Q_S(L)a)_n=\sum_{j=1}^{J} m_j\,Q_S(\mu_j)\,\mu_j^{\,n}\; }.
\]
因此对所有 \(\mu_j\in S\) 的模态，其贡献被 \textbf{严格消去}（因为 \(Q_S(\mu_j)=0\)）。
\end{proposition}
\begin{proof}
对单项 \(\mu^n\) 有 \(L(\mu^n)=\mu^{n+1}=\mu\cdot \mu^n\)，从而对任意多项式 \(Q\) 有 \(Q(L)(\mu^n)=Q(\mu)\mu^n\)。
对 \(\sum_j m_j\mu_j^n\) 线性叠加即得结论。证毕。
\end{proof}

\begin{corollary}[真实输入 40 状态核的整数谱消去滤波器]\label{cor:real-input-40-spectral-annihilation-filter}
\leanverified{paper\_real\_input\_40\_spectral\_annihilation\_filter}
令 \(a_n=\Tr(M^n)\)。
取整数多项式
\[
Q_{\mathbb{Z}}(x):=(x^2-3x+1)(x^2+x-1)=x^4-2x^3-3x^2+4x-1,
\]
并定义滤波后的序列
\[
c_n:=(Q_{\mathbb{Z}}(L)a)_n
=a_{n+4}-2a_{n+3}-3a_{n+2}+4a_{n+1}-a_n.
\]
则对所有 \(n\ge 1\) 有完全显式闭式
\[
\boxed{\;c_n=-2-15(-1)^n\;}
\]
（即 \(c_n\) 在 \(-17\) 与 \(13\) 之间交替）。
\end{corollary}
\begin{proof}
由命题 \ref{prop:finite-rh-40} 的谱分解，\(M\) 的非零特征值为
\(\{\lambda,\ -\sqrt{\lambda},\ 1^{(2)},\ (-1)^{(3)},\ \lambda^{-1},\ \lambda^{-1/2}\}\)。
其中
\(\lambda,\lambda^{-1}\)
是 \(x^2-3x+1\) 的两根，而
\(-\sqrt{\lambda},\lambda^{-1/2}\)
是 \(x^2+x-1\) 的两根，故 \(Q_{\mathbb{Z}}\) 严格消去这四个模态。
因此 \(c_n\) 只由 \(\pm 1\) 的贡献决定：
\[
c_n=2\,Q_{\mathbb{Z}}(1)\cdot 1^n+3\,Q_{\mathbb{Z}}(-1)\cdot (-1)^n.
\]
直接代入 \(Q_{\mathbb{Z}}(1)=-1\)、\(Q_{\mathbb{Z}}(-1)=-5\) 得
\(
c_n=-2-15(-1)^n
\)。
证毕。
\end{proof}

\begin{remark}[二阶滤波（系数在 \texorpdfstring{$\QQ(\sqrt5)$}{Q(sqrt5)}）]\label{rem:real-input-40-spectral-annihilation-filter-q2}
在本核上，若允许系数在 \(\QQ(\sqrt5)\)，则最低阶消去主/达界两模态 \(\{\lambda,-\sqrt{\lambda}\}\) 的滤波器为
\(
Q_2(x)=(x-\lambda)(x+\sqrt{\lambda})
\)，
对应的二阶差分关系为
\[
(Q_2(L)a)_n=a_{n+2}+(\sqrt{\lambda}-\lambda)a_{n+1}-\lambda^{3/2}a_n.
\]
当 \(\lambda=\varphi^2\) 时，\(\sqrt{\lambda}-\lambda=-1\) 且 \(\lambda^{3/2}=\varphi^3\)，因此可写成
\(
(Q_2(L)a)_n=a_{n+2}-a_{n+1}-\varphi^3 a_n
\)。
\end{remark}

\begin{proposition}[平方根项的相位共振公式：临界模谱与 \texorpdfstring{$d=2$}{d=2} 项的干涉]\label{prop:finite-rh-sqrt-resonance-general}
\leanverified{paper\_finite\_rh\_sqrt\_resonance\_general}
设 \(M\) 为 primitive 非负整数矩阵，\(\lambda=\rho(M)>1\)。并假设满足 kernel RH 的谱半径界：除 Perron 根外所有特征值 \(\mu\) 均满足 \(|\mu|\le \sqrt{\lambda}\)。
记临界模谱集合
\[
\mathcal{C}:=\{\mu\in\mathrm{spec}(M):\ |\mu|=\sqrt{\lambda}\},
\qquad
\omega_\mu:=\mu/\sqrt{\lambda}\ \ (|\omega_\mu|=1),
\]
并记严格次临界半径
\[
R:=\max\{|\mu|:\ \mu\in\mathrm{spec}(M),\ |\mu|<\sqrt{\lambda}\},
\qquad
\Lambda:=\max\{R,\ \lambda^{1/3}\}.
\]
令 \(a_n:=\Tr(M^n)\)，并定义 primitive 轨道数
\[
p_n:=\frac{1}{n}\sum_{d\mid n}\mu(d)\,a_{n/d}.
\]
则当 \(n\to\infty\) 时有渐近展开
\[
\boxed{\;
p_n
=\frac{\lambda^n}{n}
+\frac{\lambda^{n/2}}{n}\Bigl(\sum_{\mu\in\mathcal{C}}\omega_\mu^{\,n}-\mathbf{1}_{2\mid n}\Bigr)
+O\!\left(\frac{\Lambda^n}{n}\right)\; }.
\]
\end{proposition}
\begin{proof}
由谱分解（Jordan 项并入严格次级项）可写
\[
a_n=\Tr(M^n)=\lambda^n+\sum_{\mu\in\mathcal{C}}\mu^n+O(R^n)
=\lambda^n+\lambda^{n/2}\sum_{\mu\in\mathcal{C}}\omega_\mu^{\,n}+O(R^n).
\]
由 Möbius 反演，
\[
np_n=\sum_{d\mid n}\mu(d)\,a_{n/d}
=a_n-\mathbf{1}_{2\mid n}a_{n/2}+\sum_{\substack{d\mid n\\ d\ge 3}}\mu(d)\,a_{n/d}.
\]
其中对 \(d\ge 3\)，主项 \(\lambda^{n/d}\) 至多为 \(\lambda^{n/3}\)，并且余项 \(O(R^{n/d})\) 至多为 \(O(R^{n/3})\)，因此
\[
\sum_{\substack{d\mid n\\ d\ge 3}}\mu(d)\,a_{n/d}=O(\Lambda^n).
\]
另一方面，
\[
a_{n/2}=\lambda^{n/2}+O(\lambda^{n/4})+O(R^{n/2}).
\]
代入并整理得到
\[
np_n=\lambda^n+\lambda^{n/2}\sum_{\mu\in\mathcal{C}}\omega_\mu^{\,n}-\mathbf{1}_{2\mid n}\lambda^{n/2}+O(\Lambda^n),
\]
两边除以 \(n\) 即得命题结论。证毕。
\end{proof}
